\begin{proof}[Beschreibung des Agda-Beweises von \texttt{ardens-left}]
Sei $\Eps \notin A$ und sei $X$ eine Lösung der Gleichung
$X=A\cdot X\cup B$.
Der Agda-Beweis zeigt die Inklusion
\[
  X \subseteq A^*B .
\]
Zusammen mit der bereits im manuellen Beweis verwendeten Inklusion
$A^*B\subseteq X$ folgt daraus die Eindeutigkeit der Lösung.

Der eigentliche Induktionsbeweis steckt in der Hilfsaussage
\texttt{ardens-left-gen}.
Gezeigt wird die stärkere Aussage:
\[
  \forall n\in\mathbb{N}:\forall w\in\Sigma^*:
  |w|<n \text{ und } w\in X
  \quad\Rightarrow\quad
  w\in A^*B .
\]
Wir führen eine Induktion über $n$ durch.

I.A. ($n=0$): Es gibt kein Wort $w$ mit $|w|<0$, also ist nichts zu zeigen.

I.S. ($n\rightsquigarrow n+1$): Sei $w\in\Sigma^*$ mit $|w|<n+1$ und
$w\in X$.
Wegen $X=A\cdot X\cup B$ gibt es zwei Möglichkeiten.
\begin{itemize}
 \item Falls $w\in B$, dann ist
   $w=\Eps\cdot w$ mit $\Eps\in A^0$.
   Also gilt $w\in A^*B$.
 \item Falls $w\in A\cdot X$, dann gibt es Wörter $u,v$ mit
   $w=uv$, $u\in A$ und $v\in X$.
   Der Fall $u=\Eps$ ist unmöglich, denn daraus würde
   $\Eps\in A$ folgen, im Widerspruch zur Annahme.
   Also ist $u\ne\Eps$.
   Dann ist $v$ ein echtes Restwort von $w$; aus $|w|<n+1$ folgt
   $|v|<n$.
   Nach Induktionsvoraussetzung gilt daher $v\in A^*B$.
   Also gibt es ein $j$ sowie Wörter $y,z$ mit
   $v=yz$, $y\in A^j$ und $z\in B$.
   Wegen $u\in A$ folgt $uy\in A^{j+1}$ und damit
   $w=uv=uyz\in A^{j+1}B\subseteq A^*B$.
\end{itemize}

Damit ist die Hilfsaussage bewiesen.
Für \texttt{ardens-left} selbst wählt der Agda-Beweis zu einem beliebigen
$w\in X$ den Wert $n=|w|+1$.
Dann gilt $|w|<n$, und die Hilfsaussage liefert sofort $w\in A^*B$.
Also gilt $X\subseteq A^*B$.
\end{proof}
